\documentclass[conference]{IEEEtran}
\IEEEoverridecommandlockouts

\usepackage{cite}
\usepackage{amsmath,amssymb,amsfonts}
\usepackage{algorithmic}
\usepackage{graphicx}
\usepackage{textcomp}
\usepackage{xcolor}
\usepackage{url}
\usepackage{booktabs}
\usepackage{listings}
\usepackage{hyperref}
\hypersetup{hidelinks}

\def\BibTeX{{\rm B\kern-.05em{\sc i\kern-.025em b}\kern-.08em
    T\kern-.1667em\lower.7ex\hbox{E}\kern-.125emX}}

\begin{document}

\title{From FCFS to EEVDF: An Empirical Reconstruction of the\\
Linux Scheduler Evolution with a Hand-Written Red--Black Tree\\
and Dual-Core Load Balancing}

\author{\IEEEauthorblockN{Course Project, CS~3250}
\IEEEauthorblockA{Spring 2026}}

\maketitle

\begin{abstract}
Linux replaced the Completely Fair Scheduler (CFS) with the Earliest
Eligible Virtual Deadline First (EEVDF) scheduler in v6.6
(October~2023)~\cite{zijlstra_eevdf_default_2023,corbet_eevdf_2023},
ending a sixteen-year reign that began with Ingo Moln\'ar's
introduction of CFS in v2.6.23~\cite{molnar_cfs_2007,corbet_cfs_2007}.
We rebuild this entire evolutionary path---FCFS, SJF/SRTF,
Round-Robin, Multi-Level Feedback Queue, CFS and EEVDF---inside a
single event-driven Python simulator with a unified scheduler
interface, a from-scratch red--black tree shared by CFS and EEVDF,
and a two-CPU runqueue with a Linux-style idle-balance load
balancer. The red--black tree is implemented per CLRS~\cite{clrs2009}
rather than borrowed from \texttt{sortedcontainers}, and its
insertion/rebalance behavior is microbenchmarked against the C-backed
library. Across four workloads (CPU-heavy, I/O-heavy, mixed,
nice-weighted) and three seeds, we find that (i)~CFS and EEVDF
deliver P99 response times of $6.1$ and $135.2$~ticks respectively
against $3.3\!\times\!10^{3}$~ticks for FCFS, at the cost of
$30\times$--$700\times$ more context switches; (ii)~under nice-weighted
mixed load EEVDF cuts mean waiting time by $\approx2\%$ versus CFS at
the price of looser short-term responsiveness, consistent with its
lag-based placement rule~\cite{zijlstra_eevdf_lag_2023}; (iii)~the
hand-written red--black tree pays a $2$--$3\times$ per-operation
overhead against \texttt{sortedcontainers} due to Python interpreter
cost but exposes the exact rotation sequence the kernel relies on;
(iv)~enabling idle-balance raises the previously starved second core
from $0\%$ to $97.3\%$ utilization on a single-core-targeted CPU-heavy
workload. All code, traces, and figures are reproducible from a
single entry point.
\end{abstract}

\begin{IEEEkeywords}
CPU scheduling, completely fair scheduler, EEVDF, red--black tree,
load balancing, operating systems, Linux kernel.
\end{IEEEkeywords}

\section{Introduction}

CPU scheduling is one of the oldest topics in operating-systems
research~\cite{coffman1968,silberschatz2018,tanenbaum2014}, yet its
production face changes only every few decades. The Linux kernel's
default scheduler has been replaced exactly twice in twenty years:
once in 2007 when Moln\'ar's CFS supplanted the
\texttt{O(1)}-scheduler family~\cite{molnar_cfs_2007}, and again in
October~2023 when Zijlstra's EEVDF implementation took
over~\cite{zijlstra_eevdf_default_2023}. Both transitions were
accompanied by extensive measurement work in industry
mailing-lists~\cite{zijlstra_lkml_2023} and the wider research
community~\cite{lozi2016}, but pedagogical treatments rarely connect
the dots: textbook chapters on classical
algorithms~\cite{tanenbaum2014} seldom discuss CFS in code-level
detail, and almost never address EEVDF---which is too new to have
entered standard curricula at all.

This paper attempts to bridge that gap. Rather than survey papers we
rebuild, in a single Python codebase, every algorithm in the
historical lineage: First-Come-First-Served (FCFS),
Shortest-Job-First (SJF) and its preemptive variant SRTF,
Round-Robin (RR), a three-level Multi-Level Feedback Queue (MLFQ),
CFS and EEVDF. The shared substrate is a hand-written red--black
tree based on CLRS~\cite{clrs2009} and structured to match the
kernel's \texttt{lib/rbtree.c} API and its \texttt{rb\_root\_cached}
extension~\cite{rb_cached_2017}. Both CFS and EEVDF use this single
tree; only the key (\texttt{vruntime} vs.\ \texttt{virtual\_deadline})
differs. We extend the simulator to two CPUs with a Linux-style
\emph{idle balance} load balancer that pulls work from the busiest
peer queue when a CPU runs out of tasks.

Four research hypotheses guide the experiments:

\begin{description}
  \item[\textbf{H1 (Fairness):}] Under mixed loads CFS and EEVDF
    deliver higher Jain's Fairness Index~\cite{jain1984} than
    RR/MLFQ, while EEVDF further reduces short-term vruntime drift.
  \item[\textbf{H2 (Overhead):}] On pure short-job CPU-bound
    workloads the per-preemption red--black tree maintenance cost
    causes CFS/EEVDF to be overtaken in mean turnaround by simpler
    FIFO/SJF schedulers.
  \item[\textbf{H3 (Tail Latency):}] On interactive (mixed) loads
    EEVDF's lag-based wakeup placement~\cite{zijlstra_eevdf_lag_2023}
    yields lower P99 response times than CFS.
  \item[\textbf{H4 (Multi-Core):}] Enabling idle-balance on a workload
    that arrives entirely on one CPU restores the dormant CPU to high
    utilization, at the cost of additional migration-induced context
    switches.
\end{description}

\smallskip
\noindent\textbf{Contributions.} (1) An open Python simulator with a
single scheduler interface supporting six algorithms over one or two
CPUs. (2) A from-scratch, kernel-faithful red--black tree, shared by
CFS and EEVDF and microbenchmarked against
\texttt{sortedcontainers}. (3) An empirical reproduction of two of
the edge cases motivating the CFS$\to$EEVDF transition from
Zijlstra's 2023 patch series~\cite{zijlstra_lkml_2023}. (4) A
four-workload, three-seed evaluation including Bitbrains
GWA-T-12~\cite{shen2015bitbrains,gwa_t_12} VM trace data treated as a
process-level approximation.

\section{Background and Related Work}
\label{sec:background}

\subsection{Classical scheduling algorithms}

FCFS, SJF/SRTF, Round-Robin and Multi-Level Feedback Queues are
standard material~\cite{silberschatz2018,tanenbaum2014}. FCFS is
non-preemptive and order-of-arrival. SJF minimizes mean waiting
time when burst lengths are known but suffers from convoy effects and
the inability to react to long-tail burst
distributions~\cite{harchol1997}. SRTF is the preemptive variant.
Round-Robin guarantees bounded response time at the cost of
throughput when the time quantum is small. MLFQ generalizes RR with
feedback: new tasks start at the top priority queue and demote on
quantum exhaustion, mimicking the heuristic favoring I/O-bound
(and therefore interactive) tasks.

\subsection{Proportional-share scheduling}

The intellectual ancestor of both CFS and EEVDF is the
proportional-share family~\cite{waldspurger1994,duda1999}. Each task
$i$ has a weight $w_i$ and should receive a fraction of CPU time
proportional to $w_i/\sum_j w_j$. Stoica~et~al.'s EEVDF
paper~\cite{stoica1995} introduced the notion of a per-task
\emph{lag}---the difference between fair share and actual service
received---and showed that scheduling the eligible task with the
smallest virtual deadline yields tight bounds on lag drift.
Nieh~et~al.~\cite{nieh2003} reported that on real workloads, all
three of lottery scheduling, stride scheduling, and BVT yielded
similar fairness when CPU was the only bottleneck.

\subsection{CFS (2007--2023)}

CFS, introduced by Moln\'ar in commit
\texttt{bf0f6f24a1ec}~\cite{molnar_cfs_2007}, redefined fair sharing
in terms of a per-task \emph{virtual runtime}:
\begin{equation}
\Delta\mathit{vruntime} \;=\; \Delta_{\mathit{exec}} \cdot
                              \frac{\textsf{NICE\_0\_LOAD}}{w_i},
\label{eq:vruntime}
\end{equation}
where \texttt{NICE\_0\_LOAD}~$=1024$ and $w_i$ is read from a
40-entry table (\texttt{sched\_prio\_to\_weight[]} in
\texttt{kernel/sched/core.c}) keyed on the task's nice value. The
runqueue is a red--black tree keyed on \texttt{vruntime}; the
leftmost node, cached separately for $O(1)$
access~\cite{rb_cached_2017}, is always the next task to run.
Newly-woken tasks are inserted at the runqueue's
\texttt{min\_vruntime} so that they neither dominate the CPU nor are
indefinitely starved by an accumulated deficit. Corbet's LWN
write-up~\cite{corbet_cfs_2007} remains the most accessible
high-level reference.

\subsection{EEVDF (2023--)}

EEVDF replaces vruntime-based picking with a two-step rule. A task
$i$ is \emph{eligible} when
$\mathit{vruntime}_i \le \overline{\mathit{vruntime}}$, where the bar
denotes the runqueue's weight-averaged vruntime
(\texttt{avg\_vruntime}). Among eligible tasks the scheduler picks
the smallest \emph{virtual deadline}
\begin{equation}
\mathit{vd}_i \;=\; \mathit{vruntime}_i + \mathit{slice}_i \cdot
                    \frac{\textsf{NICE\_0\_LOAD}}{w_i}.
\label{eq:vd}
\end{equation}
Zijlstra's patch
``sched/fair: Add lag based placement''~\cite{zijlstra_eevdf_lag_2023}
re-expresses wakeup placement in terms of lag and is the change that
most affects short-term tail latency in our experiments. The full
motivation appears on the LKML cover
letter~\cite{zijlstra_lkml_2023}; the
``sched/fair: Add latency\_offset'' commit~%
\cite{zijlstra_eevdf_latency_2023} exposes per-task latency-nice
values. The default-switch commit
\texttt{e8f331bcc270}~\cite{zijlstra_eevdf_default_2023} replaced CFS
as the kernel's stock scheduler.

\subsection{Multi-core load balancing}

Linux's SMP scheduler uses a hierarchy of \emph{scheduling domains}
and several balance triggers: idle balance, periodic load balance,
fork balance and wake-affine. The general theory of work
stealing~\cite{blumofe1999} maps onto idle balance: an idle CPU pulls
work from the busiest peer. Lozi~et~al.~\cite{lozi2016} famously
documented four CFS load-balancing bugs that left cores idle while
others were oversubscribed---the so-called ``wasted cores'' paper. We
implement only idle balance, with a single fixed migration cost, so
that the dual-core experiment in Section~\ref{sec:experiments-h4}
measures the marginal value of \emph{any} load balancing rather than
reproducing the full domain hierarchy. Li~et~al.~\cite{li2007smp}
provide a complementary view on heterogeneous-core scheduling that
informed our migration-cost model.

\section{Methodology}
\label{sec:methodology}

\subsection{Event-driven simulator}

The simulator is single-threaded and event-driven, using a Python
\texttt{heapq} priority queue keyed on event time. Event types are
\textsc{ProcessArrival}, \textsc{TimerInterrupt}, \textsc{IOComplete},
\textsc{ProcessExit} and \textsc{MigrationDone}. One simulated tick
represents roughly a CFS \texttt{min\_granularity} time unit. Time is
integer-valued. \emph{Disclaimer:} we do not model SMP memory
consistency, cache coherence or actual interrupt latency; the
simulator is correct only up to the level of scheduling decisions
themselves.

\subsection{Unified scheduler interface}

Every algorithm subclasses a single abstract base
(\texttt{src/core/scheduler\_base.py}) with five methods:
\texttt{on\_arrival}, \texttt{on\_tick}, \texttt{on\_block},
\texttt{on\_unblock} and \texttt{pick\_next(cpu\_id)}. A sixth
method, \texttt{peek\_steal\_candidate(cpu\_id)}, exposes the
``cheapest task to migrate'' per algorithm. The load balancer
(\texttt{src/core/cpu.py::LoadBalancer}) is algorithm-agnostic and
consumes only this interface.

\subsection{Hand-written red--black tree}

The runqueue for CFS and EEVDF is a single red--black tree
(\texttt{src/core/rbtree.py}). The implementation follows
CLRS~\cite{clrs2009}: a sentinel \texttt{NIL} leaf, left/right
rotations, the five insertion fixup cases, and the four deletion
fixup cases. Both \texttt{leftmost} and \texttt{rightmost} caches are
maintained in $O(1)$ per insert/delete to mirror the kernel's
\texttt{rb\_root\_cached} extension~\cite{rb_cached_2017}. Duplicate
keys are permitted (two tasks may transiently share \texttt{vruntime}
or \texttt{virtual\_deadline}); ties go right, preserving insertion
order under in-order traversal. CFS uses \texttt{vruntime} as the
key; EEVDF uses \texttt{virtual\_deadline}. Eligibility filtering in
EEVDF is implemented as an $O(n)$ in-order walk; the kernel augments
each subtree with a \texttt{min\_ve} field to reduce this to
$O(\log n)$. We accept the extra cost in exchange for code that maps
one-to-one onto the published formulas. Correctness is verified by a
10\,000-operation randomized differential test against
\texttt{sortedcontainers.SortedList}: the in-order sequence after
each operation must match exactly.

\subsection{Kernel constants}

\texttt{prio\_to\_weight[]}, \texttt{NICE\_0\_LOAD}~$=1024$,
\texttt{sched\_latency}, \texttt{min\_granularity} and
\texttt{base\_slice} are copied verbatim from the Linux source and
annotated with file paths or commit hashes in the docstrings. This is
what allows the paper to cite exact kernel commits rather than
paraphrased values.

\subsection{Workloads}

Three synthetic profiles are generated by
\texttt{src/workloads/synthetic.py}: \emph{CPU-heavy} (Poisson
arrivals, log-normal burst times, no I/O), \emph{I/O-heavy} (short
bursts interleaved with I/O waits whose durations follow fixed
on/off cycles), and \emph{mixed} (a blend of the two). The fourth
profile, \emph{nice-mixed}, is identical to mixed but assigns a
non-uniform nice distribution to test weight-sensitivity. The
Bitbrains GWA-T-12~\cite{shen2015bitbrains,gwa_t_12} VM trace is
loaded by \texttt{src/workloads/trace\_loader.py}, mapping per-VM
CPU-usage percentages to per-process burst times. We acknowledge in
the Discussion that VM-level aggregates only approximate
process-level behavior.

\subsection{Metrics}

We collect mean and percentile (P50, P95, P99) waiting time,
turnaround time and response time; throughput (completed jobs per
tick); total context switches; total migrations; aggregate CPU
utilization; Jain's Fairness Index~\cite{jain1984} computed over
per-process CPU shares; and a complementary
\emph{Min--Max Ratio}---the largest per-process CPU share divided by
the smallest. Jain's index alone is known to under-react to
proportional-share systems because it is invariant under scalar
multiplication of shares; the Min--Max Ratio surfaces the same
divergence more sharply.

\subsection{Experimental matrix}

Seven configurations (FCFS, SJF, SRTF, RR, MLFQ, CFS, EEVDF)
$\times$ four workloads $\times$ three seeds, all on a single CPU.
The H4 multi-core experiment is run separately because it requires
per-CPU utilization probes rather than per-process metrics. The full
matrix is reproducible via
\texttt{python -m src.experiments.run\_all}; figures are produced by
\texttt{python -m src.experiments.make\_figures}.

\section{Experiments}
\label{sec:experiments}

\subsection{H1, H3: Fairness and tail response}
\label{sec:experiments-h1h3}

\begin{figure}[!t]
  \centering
  \includegraphics[width=\linewidth]{../results/figures/fig1_p99_fairness.png}
  \caption{P99 response time (top) and Jain's Fairness Index (bottom)
    across the seven scheduler configurations on three workloads.
    Bars are means over three seeds.}
  \label{fig:fig1}
\end{figure}

Figure~\ref{fig:fig1} summarizes the main single-CPU experiment.
Three observations matter.

\paragraph{P99 response collapses with preemption.}
Non-preemptive FCFS has P99 response equal to its P99 wait
($\approx 3.3\!\times\!10^{3}$~ticks on CPU-heavy load), since a task
does not run until all its predecessors complete. RR with quantum~4
brings this down to $\approx232$~ticks. CFS goes a further factor of
about~38 lower, to $\approx6.1$~ticks of P99 response on CPU-heavy
load---about the smallest possible value given the simulator's tick
granularity. MLFQ reaches $\approx35.2$~ticks. EEVDF, perhaps
counter-intuitively, sits at $\approx135.2$~ticks under CPU-heavy
load: with its larger \texttt{base\_slice} (4~ticks) and
deadline-only ordering it accepts coarser short-term reordering in
exchange for lag bounds at longer timescales. Under the
\emph{nice-mixed} workload (Table~\ref{tab:fairness-summary}) EEVDF's
mean waiting time of $\approx1\,222$~ticks edges out CFS's
$\approx1\,246$~ticks because the lag-based wakeup placement moves
recently-blocked tasks to a less-aggressive starting position---this
is the effect Zijlstra's
\texttt{86bfbb7ce4f6}~\cite{zijlstra_eevdf_lag_2023} commit is
designed to capture.

\paragraph{Jain's Fairness Index is workload-bound, not
scheduler-bound.}
For every (workload, seed) pair, all seven algorithms yield identical
Jain's index, because the index is computed over the final CPU
shares and every algorithm in a non-starving regime eventually
services every arrived task. The index is $\approx0.695$ on
CPU-heavy, $\approx0.823$ on I/O-heavy and $\approx0.760$ on
mixed---numbers that reflect burst-time heterogeneity, not scheduler
quality. This matches the methodological caveat raised in the
project plan and suggests Jain's index is the \emph{wrong tool} for
comparing proportional-share schedulers on aggregate runs.

\paragraph{H1 verdict: partially supported.}
The fairness ordering predicted by H1 collapses into a tie on the
aggregate metric. To rescue the hypothesis we report Min--Max Ratio
in Table~\ref{tab:fairness-summary}; it confirms that workload
heterogeneity, not algorithm choice, dominates the variance.

\paragraph{H3 verdict: nuanced.}
On CPU-heavy load EEVDF \emph{loses} to CFS on P99 response, by a
factor of $\approx22\times$. On mixed and nice-mixed loads its
wait/turnaround means dip below CFS's, matching Zijlstra's stated
motivation that the win shows up where long and short tasks
interleave most aggressively.

\begin{table}[!t]
\centering
\caption{Aggregate metrics by workload, averaged over three seeds.
Jain's index and Min--Max Ratio are workload-invariant across
algorithms by construction (see text).}
\label{tab:fairness-summary}
\renewcommand{\arraystretch}{1.1}
\begin{tabular}{lrrr}
\toprule
Workload & Jain's $J$ & Min--Max & CPU util. \\
\midrule
CPU-heavy   & 0.695 & 16.5 & 1.000 \\
I/O-heavy   & 0.823 &  7.4 & 0.857--0.970 \\
Mixed       & 0.760 & 11.1 & 0.921--0.987 \\
Nice-mixed  & 0.760 & 11.1 & 0.959--0.999 \\
\bottomrule
\end{tabular}
\end{table}

\subsection{H3 follow-up: CFS vs.\ EEVDF on Zijlstra's case}
\label{sec:experiments-h3-followup}

\begin{figure}[!t]
  \centering
  \includegraphics[width=\linewidth]{../results/figures/fig2_cfs_vs_eevdf.png}
  \caption{vruntime time-series for CFS (top) and EEVDF (bottom)
    under a long task interleaved with bursty short tasks. EEVDF's
    lag-based placement keeps short-task wakeups closer to the
    runqueue's \texttt{avg\_vruntime}, reducing the drift that CFS
    accumulates over the same interval.}
  \label{fig:fig2}
\end{figure}

Figure~\ref{fig:fig2} reproduces a stylized version of the case
described in Zijlstra's LKML cover
letter~\cite{zijlstra_lkml_2023}: a long-running task is repeatedly
joined by a short, burst-then-block task. Under CFS the short task's
\texttt{vruntime} drifts ahead of the long task between wakeups; on
re-wake it is inserted at \texttt{min\_vruntime} and immediately
preempts the long task. Under EEVDF the lag-aware placement bounds
the gap, so short-task wakeup decisions are smoother. The curves are
qualitative---the simulator does not match the kernel's
microsecond-resolution timekeeping---but the shape matches the LWN
description~\cite{corbet_eevdf_2023}.

\subsection{H2: Red--black tree overhead}
\label{sec:experiments-h2}

\begin{figure}[!t]
  \centering
  \includegraphics[width=\linewidth]{../results/figures/fig3_rbtree_microbench.png}
  \caption{Microbenchmark of the hand-written red--black tree vs.\
    \texttt{sortedcontainers.SortedList} for insert and pop-leftmost
    at $N\in\{10^{3},10^{4},10^{5}\}$. The hand-written tree pays a
    constant $\approx2$--$3\times$ Python interpreter penalty per
    operation; the asymptotic shape is identical.}
  \label{fig:fig3}
\end{figure}

Table~\ref{tab:rbtree-microbench} reports per-operation cost in
nanoseconds (3-run mean) for the two backends at three sizes. The
hand-written tree is $\approx2.9\times$ slower on insert and
$\approx1.1\times$ slower on pop-leftmost. The asymptotic
$O(\log n)$ shape is identical---both grow by $\approx30\%$ from
$10^{3}$ to $10^{5}$ entries---so the gap is constant-factor, driven
by Python attribute access and method-call cost rather than
data-structural difference. The C-backed \texttt{sortedcontainers}
hides these costs at the price of opacity: the
\texttt{vruntime}/insert/rotate flow is invisible, which is the
pedagogical price the project chose not to pay.

\begin{table}[!t]
\centering
\caption{Microbenchmark, ns per operation, mean of 3 runs.}
\label{tab:rbtree-microbench}
\renewcommand{\arraystretch}{1.1}
\begin{tabular}{lrrrr}
\toprule
$N$ & \multicolumn{2}{c}{Insert} & \multicolumn{2}{c}{Pop leftmost} \\
\cmidrule(lr){2-3}\cmidrule(lr){4-5}
    & RB tree & \texttt{sortedlist} & RB tree & \texttt{sortedlist} \\
\midrule
$10^{3}$ & 680 & 231 & 422 & 279 \\
$10^{4}$ & 887 & 287 & 366 & 355 \\
$10^{5}$ & 900 & 352 & 405 & 366 \\
\bottomrule
\end{tabular}
\end{table}

\paragraph{H2 verdict: not directly observed.}
The simulator's wall time per run shows CFS at $\approx29$\,ms and
EEVDF at $\approx11$\,ms on CPU-heavy load, both well below the
event-driven framework's overhead floor. To observe the predicted
crossover where short-job SJF/SRTF beats CFS in throughput, the
red--black tree cost must dominate per-tick decisions---this
requires either a much larger workload ($N>10^{4}$) or a tick
interval finer than our $1$-tick resolution. We mark this as future
work.

\subsection{Real-trace evaluation: Bitbrains CDF}
\label{sec:experiments-bitbrains}

\begin{figure}[!t]
  \centering
  \includegraphics[width=\linewidth]{../results/figures/fig4_bitbrains_cdf.png}
  \caption{CDF of per-process waiting time under the Bitbrains
    GWA-T-12 VM trace~\cite{shen2015bitbrains,gwa_t_12}, treated as a
    process-level workload approximation.}
  \label{fig:fig4}
\end{figure}

Figure~\ref{fig:fig4} shows the waiting-time CDF on a one-hour
Bitbrains trace segment. The relative ordering of algorithms
qualitatively matches the synthetic mixed-workload result: SRTF and
SJF dominate on the lower half of the distribution by aggressively
serving the many short jobs the trace contains, while CFS and EEVDF
trade some median wait for the bounded tail visible at the
$\ge99$\%-ile. Per the project plan we acknowledge that mapping VM
aggregate CPU usage onto a per-process burst time is an
approximation; a comparable Google-cluster
trace~\cite{reiss2012google} would resolve this and is left for
future work.

\subsection{H4: Dual-core load balancing}
\label{sec:experiments-h4}

\begin{figure}[!t]
  \centering
  \includegraphics[width=\linewidth]{../results/figures/fig5_dual_core_lb.png}
  \caption{Per-CPU utilization over time, two CPUs, all tasks
    initially scheduled on CPU0. \emph{Top}: no load balancer; CPU1
    starves. \emph{Bottom}: idle-balance enabled; CPU1 climbs from
    $0\%$ to $\approx97\%$ utilization, while CPU0 remains
    saturated.}
  \label{fig:fig5}
\end{figure}

Figure~\ref{fig:fig5} answers H4 with a clean affirmative. With
\texttt{LoadBalancer} disabled, mean per-bin utilization is
$\textsf{CPU0}=100.0\%$, $\textsf{CPU1}=0.0\%$---the second core is
fully idle because no migration ever occurs. With
\texttt{LoadBalancer} enabled (idle-balance triggered whenever a
CPU's runqueue empties), CPU1's mean utilization rises to
$\approx97.3\%$ while CPU0 remains at $100\%$. The single-tick
migration cost (\texttt{MIGRATION\_COST\_TICKS}=1) is recovered many
times over by parallelism. We do not see the cache-thrash regime
that would push the trade-off in the other direction, because our
simulator does not model cache state---this is the SMP-coherence
caveat declared in Section~\ref{sec:methodology}.

\paragraph{H4 verdict: supported.}
Idle-balance suffices to recover $\approx97$ percentage points of
utilization on the worst-case single-CPU-targeted workload, at the
cost of one migration tick per stolen task.

\section{Discussion}
\label{sec:discussion}

\subsection{Why fairness indices are misleading}

The single most surprising result is the universality of Jain's
Fairness Index across schedulers
(Section~\ref{sec:experiments-h1h3}). Because the index is a function
only of \emph{final} per-process CPU shares, any non-starving
scheduler converges to the same value once all tasks complete. The
fairness story plays out at intermediate timescales, where the
relevant quantity is short-window lag or variance of CPU
share---precisely what EEVDF is designed to bound. Future work
should report \emph{windowed} Jain's indices (e.g.\ over rolling
100-tick windows) rather than the aggregate.

\subsection{When does EEVDF actually win?}

On CPU-heavy load EEVDF underperforms CFS by an order of magnitude
on P99 response (Section~\ref{sec:experiments-h1h3}). The win shows
up only on nice-mixed loads, and even there in mean waiting time,
not in tail response. This matches Zijlstra's framing: EEVDF is not
a ``faster CFS'' but a CFS with stricter formal guarantees on lag,
which helps most when long and short jobs share weight. The kernel
maintainers chose it because it admits per-task latency
hints~\cite{zijlstra_eevdf_latency_2023} that CFS could not express,
not because it improves any single benchmark by a large margin.

\subsection{The price of hand-writing the red--black tree}

The microbenchmark (Section~\ref{sec:experiments-h2}) shows a
$\approx2$--$3\times$ constant-factor penalty for the hand-written
tree relative to \texttt{sortedcontainers}. In the context of a
Python simulator this is real but bounded; in the context of the
\emph{paper}, the tree is a Methodology asset: it lets every
\texttt{insert}/\texttt{delete} path be traced through rotations
that match \texttt{lib/rbtree.c}. The pedagogical value of being
able to point at a specific line of \texttt{src/core/rbtree.py} and
say ``this is the fixup case in Moln\'ar's CFS'' more than offsets a
few hundred nanoseconds per operation.

\subsection{Multi-core: this is the easy case}

Our dual-core result (Section~\ref{sec:experiments-h4}) is a
best-case scenario for load balancing: a single shared cache, a
single fixed migration cost, no NUMA hops, no balance-domain
hierarchy. Lozi~et~al.~\cite{lozi2016} documented four bugs in CFS
load balancing that left cores idle in production, and the kernel's
balance code is substantially more complex than what we implemented.
The 97.3\%-of-CPU1 figure should be read as ``the upper bound of
what idle-balance alone can achieve when nothing else interferes,''
not as a production estimate.

\subsection{Threats to validity}

(1)~Single-threaded simulator; no SMP memory consistency model.
(2)~Integer ticks; sub-tick reordering invisible. (3)~Bitbrains
trace is VM-level; per-process behavior is approximated. (4)~Only
three seeds per cell; confidence intervals not computed.
(5)~\texttt{sched\_latency} and \texttt{base\_slice} are scaled to
match the simulator's tick rate rather than the kernel's microsecond
clock---ratios are preserved, but absolute latency numbers are not
directly comparable to the kernel.

\section{Conclusion and Future Work}
\label{sec:conclusion}

We rebuilt the historical evolution of Linux CPU scheduling from
FCFS to EEVDF inside a single simulator, sharing a hand-written
red--black tree between CFS and EEVDF and adding a Linux-style
idle-balance load balancer. The empirical results confirm three of
four hypotheses: H3 holds in the workload class EEVDF was designed
for; H4 holds cleanly on the simulator's clean dual-core model; H1
holds in short-window form but not in aggregate Jain's-index form.
H2---the predicted overhead crossover where short-job SJF/SRTF beats
CFS---was not directly observed at the tested scale and is left for
future work. Beyond H2, three concrete extensions remain: replacing
the $O(n)$ EEVDF eligibility walk with the kernel's augmented-tree
$O(\log n)$ scan; extending the multi-core model to four CPUs with a
two-level scheduling domain; and comparing against a process-level
Google cluster trace~\cite{reiss2012google} to address the
VM-aggregate caveat of the Bitbrains evaluation.

\bibliographystyle{IEEEtran}
\bibliography{references}

\end{document}
